\section{Bài toán, mục tiêu và phạm vi}
\subsection{Bài toán nghiên cứu}
Data drift là thay đổi phân phối đầu vào $P(X)$ giữa reference và production. Cảnh báo drift không tự chứng minh mô hình suy giảm: thay đổi có thể xảy ra ở vùng mô hình vẫn dự đoán tốt. Ngược lại, sai số có thể tăng khi quan hệ $P(Y\mid X)$ thay đổi mà detector đầu vào không phát hiện. Do đó, giám sát drift và kiểm chứng chất lượng là hai nguồn bằng chứng bổ sung.

Ở thời điểm quyết định $t$, hệ thống quan sát batch production, prediction và xác suất của mô hình đang phục vụ; chỉ được đọc nhãn đã nhận hoặc được yêu cầu trong ngân sách. Cần lựa chọn \texttt{KEEP}, \texttt{REQUEST\_LABELS}, \texttt{HOLD} hoặc \texttt{RETRAIN}. Khi không đủ bằng chứng, kết quả là chưa xác định; không gán mặc định mô hình khỏe. Ước lượng không nhãn phụ thuộc giả định về shift và calibration, không bảo đảm nhận diện mọi suy giảm~\cite{atc,cbpe}.

\subsection{Câu hỏi nghiên cứu và đóng góp dự kiến}
\textbf{Định vị MLOps:} Đối tượng nghiên cứu là vòng điều phối từ giám sát đến cập nhật mô hình, không phải thuật toán ước lượng chất lượng mới. Một trigger đúng về thống kê vẫn có thể dẫn đến kết quả vận hành sai nếu dùng nhầm snapshot, chi nhãn hai lần khi retry hoặc áp dụng candidate cho champion đã thay đổi. Vì vậy, đề tài đánh giá cả hiệu quả quyết định và tính đúng đắn của quá trình thực thi.
\begin{enumerate}
  \item \textbf{RQ1 -- Hiệu quả vận hành:} Với cùng giới hạn ngân sách và recipe, kiểm chứng trước retraining đạt đánh đổi chất lượng, chi phí nhãn, số lần huấn luyện và bỏ sót suy giảm như thế nào so với trigger theo drift, estimator hoặc audit theo lịch cố định?
  \item \textbf{RQ2 -- Độ tin cậy điều phối:} Vòng Continuous Training trên Argo Workflows/Kubeflow có bảo toàn ngân sách, tính nhất quán dữ liệu--phiên bản và điều kiện cập nhật khi có sự kiện lặp, lỗi worker/pod hoặc callback đến muộn không; cần bao lâu để phục hồi sau khi thành phần lỗi hoạt động trở lại?
\end{enumerate}
Đóng góp gồm: (i) thiết kế và hiện thực vòng điều phối có kiểm chứng, trạng thái bền vững, quản lý ngân sách và truy vết; (ii) bộ đánh giá tái lập gồm replay chất lượng--chi phí và kiểm thử lỗi vận hành. Tích hợp hệ thống là đầu ra bắt buộc, không chỉ là minh họa cho notebook. Không đề xuất detector, estimator, thuật toán lấy mẫu hay lý thuyết hệ phân tán mới. Thành công được đánh giá bằng bằng chứng thực nghiệm và phân tích giới hạn, kể cả khi policy không vượt baseline.

\subsection{Ranh giới thực hiện}
\begin{itemize}
  \item \textbf{Cốt lõi:} Hai dataset nhị phân tabular, một model family có xác suất lớp; CBPE chính, ATC đối chứng; mẫu nhãn ngẫu nhiên và một recipe full-retraining. Triển khai vòng tích hợp trên AWS/K3s với Argo Workflows, Kubeflow Training Operator và Prometheus/Grafana; Docker chỉ phục vụ kiểm thử local. Kiểm chứng retry, phục hồi và cập nhật phiên bản trên môi trường thực nghiệm cùng kiến trúc production.
  \item \textbf{Giả định nhãn:} Dataset thực nghiệm có nhãn đầy đủ ở phía evaluator. Hệ thống chỉ nhận phần được tiết lộ theo ngân sách sau mỗi batch. Nhãn đúng, chi phí một đơn vị/mẫu; phần còn lại chưa được phép đọc. Yêu cầu nhãn được mô phỏng, không phải nghiên cứu hành vi người gán nhãn.
  \item \textbf{Ngoài phạm vi:} Taxonomy nguyên nhân, causal diagnosis, học policy bằng reinforcement learning, tối ưu sampler, incremental training, pseudo-labeling, nhãn trễ/sửa nhãn phức tạp và human study. Kế thừa Terraform/Ansible/GitOps, chỉ bổ sung cấu hình cần cho đề tài; không nghiên cứu high availability đa vùng, tối ưu autoscaling, distributed training hoặc canary production; không cam kết exactly-once toàn hệ thống.
\end{itemize}
Trường hợp không lấy được bất kỳ nhãn production nào chỉ dùng kiểm tra giới hạn giám sát. Không cam kết supervised retraining trên dữ liệu không nhãn. Tên dataset, cấu hình model, mức ngân sách và ngưỡng được chọn qua pilot ở tháng 1 rồi đóng băng trước test.
